% results_introduction

\section{Introduction}
\label{sec:introduction}

Salah, the ritual prayer performed five times daily, is obligatory for
practising Muslims and timed astronomically: each of the five windows moves
with the sun and the observer's location, and a small number of further
windows (during which prayer is disallowed, or during which a travelling or
menstruating person is exempt) depend on further computation and on facts
about the user's body and circumstances. A large and growing population uses
mobile applications to mediate this obligation --- computing prayer times,
sounding the call to prayer (adhan), orienting toward Mecca (qibla), and
increasingly, logging whether a prayer was performed. This mediation makes
these applications' failures different in kind from a productivity app's. A
missed prompt is not a missed reminder; for a user relying on it, it can be a
missed prayer.

The small existing literature on these applications, and the larger
literature on the design patterns they borrow, offers three framings for
thinking about what makes them succeed or fail. The self-tracking and
gamification literature suggests that streaks and completion scores --- the
mechanics a prayer tracker borrows from fitness apps --- induce guilt and
anxiety in users who cannot keep them up. Product design more broadly treats
feature richness and interface minimalism as a trade-off point along a single
axis, and treats a richer feature set as a differentiator that wins and
retains users. And accuracy is often opposed to aesthetics, as though a
user's tolerance for a wrong prayer time and a user's tolerance for a
cluttered screen were governed by the same underlying preference. Each of
these framings is plausible, each has real evidence behind it in adjacent
domains, and each turns out to mis-locate where satisfaction with a
devotional app actually comes from.

What is missing is evidence at the scale needed to test these framings against
the whole ecosystem rather than against one application. The closest existing
work is small in scope or method: a dissertation study of prayer-app design
\citep{bellar2017ipray}, an expert heuristic comparison of two named Islamic
applications \citep{azahari2025comparative}, and single-application content
analyses of review themes \citep{noor2025halaltourism,noor2026halalsuperapps}.
None mines reviews across the ecosystem with aspect-level sentiment, and none
tests the gamification-harm hypothesis in a devotional context. We do not claim
to be the first study of Islamic mobile applications, of app-store reviews, or
of gamified self-tracking; we claim to be the first, to our knowledge, to bring
ecosystem-scale review mining to bear on this set of design questions in this
domain.

We analyse 732,194 reviews across 26 Islamic prayer-time and devotional
applications on the Google Play Store, addressing six pre-specified research
questions with a mixed-methods pipeline: aspect-based sentiment analysis
across 27 functional and experiential aspects, seeded topic modelling,
regex-based demand-signal extraction, promise-versus-delivery comparison
against store descriptions, and five statistical models under a shared
false-discovery-rate correction, corroborated by two independently drawn
qualitative coding passes over tracker-related complaints.

Three of our six research questions retire a framing from the literature and
replace it with something the data supports better. Gamifying a prayer tracker
neither helps nor harms sentiment, and guilt language is \emph{rarer}, not
commoner, inside tracker discourse; what users resent is a record they cannot
trust and a prompt that never arrives (RQ1, \S\ref{sec:rq1}). Accuracy and
interface complaints cost the same \emph{as classes}, but that average conceals
a severity gradient, and a large share of accuracy complaints arrive inside
four- and five-star reviews, invisible to a rating-based read of severity (RQ2,
\S\ref{sec:rq2}). And feature count does not predict interface-bloat complaints
in either direction; what users mean by bloat turns out to be a property of the
interface, which bloat and interface complaints co-occurring far above base
rate inside individual reviews shows directly (RQ5, \S\ref{sec:rq5}). Two further questions test whether
distinctive features confer competitive advantage (they do not; users cite core
reliability instead, RQ3, \S\ref{sec:rq3}) and whether bio-spiritual needs such
as menstrual accommodation are met (they are not, though with only three
implementing applications we treat this as an unmet-need finding rather than a
causal one, RQ4, \S\ref{sec:rq4}). A sixth synthesises these into two distinct
failures, a coverage gap and a delivery gap requiring different remedies (RQ6,
\S\ref{sec:rq6}).

This paper makes four contributions: an ecosystem-scale account of
devotional-app quality across six pre-specified research questions, with the
false-discovery-rate denominator and every non-inferable term reported rather
than dropped; a reframing of gamification harm in
religious practice as record integrity, prompting reliability, and the
application's model of worship rather than guilt; a demonstration that
devotional register breaks off-the-shelf sentiment and aspect-extraction
tooling, which we treat as a finding about applying general-purpose NLP to
religious text rather than a data-cleaning footnote; and two sampling hazards
in app-store research, each demonstrated with numbers rather than asserted ---
star ratings are an unsafe proxy for defect severity, and single-application
qualitative sampling can manufacture what look like genre-level themes.
